\section{Allineamento Esterno}
\label{sec:allineamento}

\subsection{equivalentClass verso Wikidata}

L'ontologia dichiara 4 allineamenti \texttt{owl:equivalentClass} verso entit\`a
Wikidata, riportati nella Tabella~\ref{tab:align_class}.

\begin{table}[h]
\centering
\caption{Allineamenti equivalentClass verso Wikidata}
\label{tab:align_class}
\begin{tabular}{lll}
\toprule
\textbf{Classe VG} & \textbf{Entit\`a Wikidata} & \textbf{Descrizione Wikidata} \\
\midrule
VideoGame & wd:Q7889      & Video game \\
Developer & wd:Q210167    & Video game developer \\
GameEngine& wd:Q193564    & Game engine \\
Genre     & wd:Q65956376  & Video game genre \\
\bottomrule
\end{tabular}
\end{table}

\subsection{equivalentProperty verso Dublin Core}

Tre datatype property sono allineate con il vocabolario Dublin Core \citep{dublincore}
tramite \texttt{owl:equivalentProperty} (Tabella~\ref{tab:align_prop}).

\begin{table}[h]
\centering
\caption{Allineamenti equivalentProperty verso Dublin Core}
\label{tab:align_prop}
\begin{tabular}{lll}
\toprule
\textbf{Propriet\`a VG} & \textbf{Propriet\`a DC} & \textbf{Label DC} \\
\midrule
releaseDate    & dcterms:issued       & Date Issued \\
gameName       & dcterms:title        & Title \\
gameDescription& dcterms:description   & Description \\
\bottomrule
\end{tabular}
\end{table}

\subsection{Vantaggi dell'Allineamento}

L'allineamento esterno, noto in letteratura come \textbf{ontology alignment}
(cfr.~Lez.~13), offre tre vantaggi principali:

\begin{enumerate}[nosep]
  \item \textbf{Interoperabilit\`a}. Un'applicazione che consuma dati Dublin Core
    pu\`o automaticamente interpretare \texttt{vg:releaseDate} come sinonimo di
    \texttt{dcterms:issued}, senza bisogno di mapping manuali;
  \item \textbf{Linkage automatico}. Le entit\`a Wikidata esterne referenziate
    tra le istanze (es. \texttt{wd:Q12345}) sono automaticamente riconosciute come
    istanze di \texttt{vg:VideoGame} se Wikidata le classifica come Q7889;
  \item \textbf{Riuso di standard consolidati}. Anzich\'e reinventare vocabolari,
    l'ontologia si appoggia a standard ampiamente adottati (Dublin Core per
    metadati bibliografici, Wikidata per identificazione entit\`a), seguendo
    il principio \textit{Linked Data} di riusare URI esistenti.
\end{enumerate}

\subsection{Limitazione Attuale}

L'allineamento \`e dichiarato a livello dello schema ma non viene sfruttato dal
reasoner OWL 2 RL, poich\'e \texttt{owlrl} opera solo sul namespace locale.
Per sfruttare appieno questi assiomi, sarebbe necessario un reasoner DL completo
(HermiT, Pellet) che possa importare anche le definizioni di classe da Wikidata.
